\section{Introduction}
\label{sec:introduction}

Robot soccer is a fast, real-world robotics domain where robust self-localization is an important prerequisite for useful play. A robot must therefore continuously estimate where it is on the field, with centimeter-level precision, from partial, intermittent, robot-perspective observations under a limited field of view while moving at high speed.

This work takes inspiration from biological navigation, where animals localize by combining self-motion cues with external spatial cues such as landmarks~\cite{tang_learning_2024},
suggesting a robot localization model in which motion input and landmark observations jointly maintain a latent belief about position on the field.
The generalized-coordinate formulation of active inference naturally integrates motion cues, such as velocity and acceleration, with landmark observations
in a shared belief over field pose. It thereby turns the biological intuition into a compact prediction-error minimization algorithm.

Active inference has been used as a framework for self-localization and navigation, but existing work does not provide a compact real-time localization formulation tailored to the fast-paced, centimeter-level requirements of robot soccer. For this setting, we derive a generalized-coordinate belief update on a known soccer field that combines motion and landmark prediction errors and reduces to compact matrix calculations that show a robust working solution in a simulation environment.

The paper makes the following contributions:
\begin{itemize}
    \item it formulates robot-soccer self-localization as generalized-coordinate active-inference belief updating on a known soccer field;
    \item it shows, in simulation experiments, that the resulting method is efficient, robust, fast, and precise enough for RoboCup soccer.
\end{itemize}
